% Part IV: the unified picture

The three parts assemble into a single structure.  The Dedekind
comparison is, by \cite{PaperW}, equivalent to one homotopical
statement---chain approximation on all fibrant objects---and the two
combinatorial statements to which its type-level upgrade was localized
both feed that same core:
\[
  \small
  \boxed{\;
  \begin{array}{c}
  \textbf{master:}\ \ \Ch(X)\hookrightarrow X\ \text{a type equivalence
  for all fibrant }X\\
  \Longleftrightarrow\ \text{Dedekind comparison}\\[4pt]
  \text{(a) well-foundedness of the subatom order}\ \Longleftarrow\
  n\text{-coskeletality of atoms}\\[2pt]
  \text{(b) intrinsic chain-anchored asphericity}\ \Longleftrightarrow\
  \text{master on anchored }W\subseteq\text{atoms}
  \end{array}\;}
\]
Statement (a) is a \emph{finiteness}, orthogonal to the homotopy content:
it secures the well-founded induction that drives the assembly, and it
reduces to the $n$-coskeletality of atoms---a skeletal determinacy for
each fixed cube, untouched by the unbounded \emph{Kan} coskeletality of
non-median-fibre atoms.  The natural shortcut, a height bound on the
generation arity, is \emph{false} from $\cube^4$ on
(Proposition~\ref{prop:nobound}), which is exactly why the finiteness
must be carried by determinacy rather than by an elementary measure.
Statement (b) is not independent at all: it is the master statement
restricted to the anchored subobjects of atoms.  So the Dedekind endgame
has a single homotopical core and two auxiliaries---one a finiteness, one
an instance of the core.

Part~\ref{part:cx} explains \emph{why} the whole edifice pivots on
median-closure.  The realizability engine underlying atoms, coskeleta and
approximation is a monotone-selection constraint problem, and two
thresholds meet at the median mechanism.  For a \emph{fixed} generator,
median-closure of the fibres is exactly the condition that keeps the
canonical median a majority polymorphism, hence gives strict width two.
For the \emph{ambient} order-convex template, the bounded-width and Taylor
structure collapses at dimension three, so its $\mathrm{CSP}$ jumps to
unbounded width and $\mathrm{NP}$-completeness (Theorem~\ref{thm:cx}).  This is an algebraic dichotomy of the template,
sitting exactly at the threshold that governs the bounded-width mechanism
studied in \cite{PaperW}---the median majority polymorphism and its
strict-width/coskeletality program, whose passage from strict width two to
bounded coskeletality is itself left conjectural there; whether the
template dichotomy
transfers to a hardness statement about cubical realizability itself, or
whether (as our computations suggest) fixed-generator realizability stays
tractable, is a separate question we do not settle.  Even so, the
coincidence of the algebraic threshold with the homotopical one is the
structural reason the Dedekind site resists the elementary methods that
settle the other cube sites, and it locates what a resolution must
overcome: prove the $n$-coskeletality finiteness, and lift chain
approximation from test to type on the anchored subobjects of atoms---the
two faces, finite and homotopical, of the same median obstruction.

\begin{table}[ht]
\centering
\small
\renewcommand{\arraystretch}{1.2}
\begin{tabular}{l l p{0.34\textwidth}}
\hline
statement & status & reduces to\\
\hline
Dedekind comparison & open & the master statement \cite{PaperW}\\
master (chain approximation) & open & ---\\
(a) well-foundedness & open & $n$-coskeletality (Part~\ref{part:a}); median-fibre reduced to pair-extension\\
(b) anchored asphericity & open & master on anchored $W$ (Part~\ref{part:b})\\
median-closed realizability & \textbf{thm} & strict width two (Part~\ref{part:cx})\\
order-convex language $L_m$, $m\ge3$ & \textbf{thm} & unbounded width, $\mathrm{NP}$-complete (Part~\ref{part:cx})\\
fixed-generator realizability & \textbf{obs.} & $\mathrm P$ in all computed cases\\
\hline
\end{tabular}
\caption{The endgame at a glance: one homotopical core, two auxiliaries,
and the complexity dichotomy that pins the median-closure boundary.}
\end{table}
